\documentclass[10pt]{article}
\usepackage[margin=.78in]{geometry}
\usepackage{amsmath,amssymb,booktabs,microtype,hyperref,enumitem}
\usepackage[T1]{fontenc}\usepackage{lmodern}
\setlist{nosep,leftmargin=*}
\title{The Inception of In-Context Learning:\\Minimal Training Structure for Contextual Candidate Promotion}
\author{Paper 0.8 --- Controlled ICL Mechanism Research Line}
\date{Draft research plan --- August 2026}
\begin{document}\maketitle
\begin{abstract}
In-context learning (ICL) is often characterized through high-level analogies to regression, Bayesian inference, or implicit gradient descent. We instead ask a smaller mechanistic question: what is the minimal training distribution and minimal Transformer in which context causes a target token that is low-ranked without demonstrations to become rank one? We begin with controlled paired sequences such as \(A\!\to\!1,B\!\to\!2,C\!\to\!?\), while explicitly preventing \(C\!\to\!3\) from being learned directly in the weights. Candidate promotion is measured as a layerwise change in target rank and margin between context-free and context-conditioned inference. We progressively add only the statistical structure required for the effect---local successor relations, multiple chains, task diversity, positional variation, and episode-specific mappings---and stop scaling once the smallest competent regime is found. We then compare models immediately below and above the behavioral threshold using residual, self-attention, feed-forward, head-level, and causal replacement measurements. The goal is a bottom-up account of how contextual evidence first creates a prediction that the frozen weights alone do not favor.
\end{abstract}
\section{Question}
Let \(Y\) be a candidate with low context-free rank, \(\operatorname{rank}(Y\mid X)\gg1\). We seek the smallest training distribution \(\mathcal D^\star\) and smallest model \(M^\star\) for which \(\operatorname{rank}(Y\mid D_{\rm ctx},X)=1\), while \(Y\) remains low-ranked without \(D_{\rm ctx}\).
Define
\[
G_{\rm ICL}(Y)=m(Y\mid D_{\rm ctx},X)-m(Y\mid X),\quad
m(Y)=z(Y)-\max_{y\ne Y}z(y).
\]
The threshold is crossed when contextual margin becomes positive on held-out relations and the effect disappears under matched shuffled/wrong-context controls.

\section{Minimal motivating example}
Train local orders \(A\to B,B\to C\) and \(1\to2,2\to3\), but never train \(C\to3\) or the complete paired mapping. Test
\[
A\;1,\qquad B\;2,\qquad C\;?
\]
and ask whether context aligns the two local chains strongly enough to promote \(3\).

\section{Minimal-dataset ladder}
D0: random/no useful structure.\\
D1: isolated local successor chains only.\\
D2: many isomorphic independent chains.\\
D3: minimal paired/alignment exposure if D1--D2 fail.\\
D4: episode-specific random permutations preventing global mapping memorization.\\
D5: richer arbitrary/compositional relations only after the minimal mechanism is understood.

At each stage add exactly one structural ingredient.

\section{Dataset and model phase diagrams}
Search over number of relation families, examples per family, task/mapping entropy, chain length, symbol reuse, template variation, and position variation. Start with
\[
L\in\{1,2,3,4,6\},\qquad d_{\rm model}\in\{8,16,32,64\},
\]
and one or two heads where feasible. The goal is a behavioral phase boundary, not maximum benchmark performance.

\section{Controls}
Every positive cell compares no demonstrations, correct demonstrations, shuffled outputs, wrong-chain demonstrations, same-length irrelevant context, order perturbations, random initialization, held-out symbols, and held-out relation families. Audit that the critical direct mapping never occurs in training and that the target remains low-ranked without context.

\section{Layerwise mechanism}
At every boundary record residual state, SA update, FF update, target logit, rank, margin, top-$k$ candidates, and entropy. Compare context-free and context-conditioned trajectories. The first question is where \(Y\) becomes a serious candidate; the second is where it becomes rank one.

\section{Mechanistic hypotheses}
Test whether attention aligns corresponding relational roles, transports a successor/order state, or directly increases support for the target. Test whether FF already encodes reusable successor transformations, maps an attention-transported relation into the target direction, or simply sharpens a candidate. These are competing hypotheses.

\section{Causal experiments}
Ablate heads; zero or replace SA/FF updates; patch states from correct, shuffled, wrong, role-matched, and matched-norm controls; block individual demonstration tokens; and remove one local relation from training and retrain. Prefer causal loss of contextual promotion over attention-map interpretation.

\section{Before/after acquisition}
Find closely matched \(\mathcal D^-\) and \(\mathcal D^+\) conditions immediately below and above the ICL threshold. Compare head utility, SA/FF updates, residual target projection, relation-role similarity, cross-layer covariance, and weights. This threshold comparison is central.

\section{Minimality}
After finding a positive regime, delete one training family, relation type, positional regularity, symbol overlap, layer, head, or width. If the effect survives, simplify again.

\section{Interpretation}
Do not assume the Transformer is performing regression, Bayesian inference, or implicit gradient descent. Those may be useful behavioral approximations at larger scale. Prefer the smallest measured mechanism consistent with the evidence:
\[
\boxed{\text{learned local relations}+\text{contextual alignment/transport}+\text{candidate promotion}.}
\]

\section{Relation to Papers 0.5--0.7}
Paper 0.5 studies predictive refinement once evidence is accessible. Paper 0.6 studies learned semantic/compositional relations. Paper 0.8 asks how context combines learned relations to create a new prediction without weight updates. Paper 0.7 remains deferred until these mechanisms justify synthesis.

\section{Scaling policy}
Do not scale merely to improve performance. Scale only if the minimal model cannot express the capability, a mechanistic hypothesis requires additional components, or external validity becomes the explicit question.

\section{Expected contribution}
The target contribution is not another demonstration that Transformers perform ICL. It is an account of the \emph{inception} of ICL: the smallest statistical environment and neural architecture in which a frozen Transformer uses context to promote a prediction that its weights alone do not support, together with a causal account of where that probability mass comes from.
\end{document}
